\subsection{Naturality under an abstract isomorphism}

\begin{lemma}[One projective transformation for both coefficient lines]
\label{lem:common-g-functoriality}
Let \(\varphi:\widehat D_R\xrightarrow{\sim}\widehat D_{R'}\)
be an abstract isomorphism, with \(R,R'\in G_4^\circ\).
It induces an isomorphism \(\sigma:\Sigma_0(R)\to\Sigma_0(R')\)
and a grading-preserving isomorphism of their normal cones.  There
is one \(g\in\PGL(V)\), independent of the point of \(\Sigma_0(R)\),
which transports both recovered nonzero Pluecker component lines:
\[
 [Pw_{R'}]=[P((\bigwedge^4\Sym^2g)w_R)],\qquad
 [Qw_{R'}]=[Q((\bigwedge^4\Sym^2g)w_R)].
\]
The corresponding vanishing statements hold for zero components.
No extension of \(\varphi\) to either ambient Grassmannian is required.
\end{lemma}
\begin{proof}
The deepest stratum is intrinsic by
Lemma~\ref{lem:deepest-normal-cone}.  If \(\mathfrak b\) and
\(\mathfrak b'\) are its ideals in the two failure schemes, pullback
by \(\varphi\) identifies every power of these ideals.  Hence it
induces the graded algebra isomorphism
\begin{equation}\label{eq:graded-cone-functoriality}
 \sigma^*\!\left(\bigoplus_{m\ge0}
                  \mathfrak b'^m/\mathfrak b'^{m+1}\right)
       \simeq \bigoplus_{m\ge0}\mathfrak b^m/\mathfrak b^{m+1}.
\end{equation}
The symmetric algebra of degree one surjects canonically onto either
algebra.  Thus \eqref{eq:graded-cone-functoriality} also identifies
its homogeneous kernel, not merely the cone's reduction; this is the
functorial normal-cone construction \cite{StacksNormalCone}.
Put \(\mathcal G=\gr_{\mathfrak b}\OO_{\widehat D_R}\),
where \(\mathfrak b\) is the deepest-stratum ideal.  Its degree-one
part is the conormal bundle \(\mathcal G_1=V^*\otimes\mathcal K_0\);
its dual, not \(\mathcal G_1\) itself, is
\(\Hom(\mathcal K_0,V)\).  Thus the vector bundle is recovered
before choosing either tensor factor.

Here is the whole construction in one diagram.  The arrows
\(\rightsquigarrow\) mean assignments functorial under abstract
isomorphisms, not morphisms or chosen retractions between the displayed
schemes.  Write \(\mathcal A=\Sym\mathcal G_1\), with homogeneous
kernel \(\mathcal I\), and \(\mathcal D=\sqrt{\mathcal I}\).
\begin{equation}\label{eq:full-intrinsic-chain}
\begin{gathered}
 \widehat D_R\rightsquigarrow B=\Sigma_0(R)
 \rightsquigarrow\mathcal G\rightsquigarrow
           (\mathcal G_1,\mathcal A,\mathcal I,\mathcal D),\\
 \Spec_B(\mathcal A/\mathcal D)
 \rightsquigarrow\{\rank\le1\}
 \rightsquigarrow\operatorname{Fano}_{4}
 \rightsquigarrow\Pj(V)\times B,\\
 (\mathcal I,\mathcal D)\rightsquigarrow
 L=\mathcal D_4\rightsquigarrow
 (\mathcal I:\mathcal D)_{12}=L^*\otimes\mathcal I_{16}
 \rightsquigarrow\{\C\beta_{R,\lambda}\}_{\lambda\in\Lambda}
 \rightsquigarrow([Pw_R],[Qw_R]).
\end{gathered}
\end{equation}
The first row uses the reduced-singular-locus construction and powers
of its ideal.  The second uses reduced rank strata, the closed relative
Fano scheme of Lemma~\ref{lem:rank-one-fano-scheme}, and its uniquely
projectively trivial component.  The third uses the intrinsic colon,
Cauchy projection and contraction of the entire right factor, followed
by exterior duality.  It involves no selected generator of \(L\).
Every row is transported by \(\varphi\); comparison of the two
trivial ruling bundles, after pullback along \(\sigma\), is what
produces \([g]\).  There is no coordinate map \([g]\) attached
to an individual scheme before comparing its projective factors.

The relative rank-one Segre variety and its two ruling parameter
bundles are preserved.  Their distinction by projective triviality
was proved in Lemma~\ref{lem:deepest-normal-cone}; in particular the
\(V\)-ruling cannot be exchanged with the other ruling.  Lemma~\ref{lem:oriented-segre-descent}, applied after
pullback by \(\sigma\), gives a regular isomorphism of the
trivial \(\Pj(V)\)-bundles.  It explicitly accounts for the
line-bundle ambiguity of its local linear lifts.  Since
\(\Sigma_0(R)\) has only constant global regular functions,
the resulting \(\PGL(V)\)-valued morphism is constant.
Denote its value by \(g\).

Fix a point \(U=K_0\) and its image \(U'=K'_0\).  A linear
isomorphism preserving both rank-one rulings has, after a lift of
\(g\), the form
\(\alpha(T)=gTh^{-1}\) for some isomorphism \(h:U\to U'\).
Indeed the two Segre factors give projective linear maps.  After
lifting them, the remaining linear map fixes each decomposable line;
sums of tensors with a common factor force its scalar to be constant.
This scalar can be absorbed in \(h\).

The only frame-dependent display is the following naturality square;
all its objects were already recovered in \eqref{eq:full-intrinsic-chain}.
A change of a local lift \(g_i\) to \(u_{ij}g_i\) is a
line-bundle transition, not another projective coordinate map.  A
homogeneous degree-twelve coefficient functor acquires the common
factor \(u_{ij}^{12}\) (or its inverse on a dual).  It therefore
preserves the recovered projective lines.  The right-frame action
is shown on the right tensor factor in the square:
\[
\begin{CD}
 \mathcal F(V)^*\otimes\mathcal F(U')
      @>{\Phi_{V,U'}}>> \Sym^{12}(V^*\otimes U')\\
 @V{\mathcal F(g)^*\otimes\mathcal F(h^{-1})}VV
      @VV{\alpha^*}V\\
 \mathcal F(V)^*\otimes\mathcal F(U)
      @>{\Phi_{V,U}}>> \Sym^{12}(V^*\otimes U).
\end{CD}
\]
In fact both routes evaluate at \(T\) as
\(\beta'(\mathcal F(g)\mathcal F(T)\mathcal F(h^{-1})\xi')\).
The induced graded map transports the determinant ideal line
and the degree-sixteen cone ideal.  By
Lemma~\ref{lem:intrinsic-residual-ideal} it therefore transports
\((\mathcal I':\mathcal D')_{12}\) to
\((\mathcal I:\mathcal D)_{12}\), that is,
\((J_{R'})_{12}\) to \((J_R)_{12}\) in the chosen fibres.
This is the intrinsic residual operation, not a choice of scalar
determinant equations.  The universal diagonal-block identity gives
\[
 \C\beta_{R,\lambda}
       =\C\mathcal F(g)^*\beta_{R',\lambda}
       \quad(\lambda\in\Lambda)
\]
in the respective left factors; \(h\) acts only on the right factor
and cannot change either left line.  Exterior duality in
Lemma~\ref{lem:readout-exterior-twist} now gives the displayed
Pluecker identities.  More explicitly the last arrow is the natural map
\[
 \bigwedge^4\Sym^2V\otimes(\det\Sym^2V)^*
             \longrightarrow\bigwedge^6(\Sym^2V)^*,
 \qquad (\det\Sym^2V)^*=(\det V)^{-5}.
\]
The factor \((\det V)^{-5}\) is common to both summands, so it
disappears projectively; it cannot introduce two independent left
transformations.  Every step arose from the
isomorphism of pairs, establishing the last assertion.
\end{proof}
